% GENERATED FILE. Edit canonical lessons, then run npm run generate:books.
% canonical-source-hash: fnv1a64:54e3445ad45417f5
% canonical-lessons: AR-C267-taqlid, AR-C267-thaqafa, AR-C267-din, AR-C267-salah, AR-C267-alaqa

\chapter{Tradition, Culture, Religion, Prayer, Relationship}
\label{ch:ar-tradition-culture-religion-prayer-relationship}
\hlchaptermodality{267}

\begin{chapteropening}
\textbf{By the end of this chapter:} \emph{I can say a tradition, culture, religion, prayer and a relationship.}

Complete the last lesson of chapter 267: I can say a tradition, culture, religion, prayer and a relationship.
\end{chapteropening}

\section[\texorpdfstring{\ar{تقليد} (taqlīd)}{تقليد (taqlīd)}]{\ar{تقليد} (taqlīd) — a tradition}

\label{lesson:AR-C267-taqlid}



\begin{quote}
Before the new one: say the Arabic for a fence, then the Arabic for a courtyard.
\end{quote}

\subsection*{You'll want to know: \ar{تقليد}}
\textbf{\ar{تقليد}} — \emph{taqlīd} — \textquotedblleft{}a tradition\textquotedblright{}.

\begin{cousinweb}
\textbf{\ar{تقليد}} (\emph{taqlīd}) is built on the native Arabic root \emph{q-l-d}, \textquotedblleft{}to hang on the neck; to follow another\textquotedblright{} (Lane). \emph{Taqlīd} is the verbal noun of Form II, \textquotedblleft{}following another's lead\textquotedblright{}.
\end{cousinweb}

\begin{grammarlens}[title={What you've built}]
One more word for this chapter.
\end{grammarlens}

\subsection*{Your turn}
\begin{itemize}
\raggedright
  \item \emph{Say it:} \emph{taqlīd}
  \item \emph{Say it:} \emph{taqlīd}, once more
  \item \emph{Say it:} say \emph{taqlīd}
  \item \emph{Recall:} say the Arabic for a fence, then the Arabic for a courtyard, then say \emph{taqlīd} again
\end{itemize}

\begin{tcolorbox}[breakable,colback=teal!4,colframe=teal!35!black,title={Before you move on}]
What does \ar{تقليد} mean? (\textquotedblleft{}A tradition\textquotedblright{}.) Say it once more.
\end{tcolorbox}
\section[\texorpdfstring{\ar{ثقافة} (thaqāfa)}{ثقافة (thaqāfa)}]{\ar{ثقافة} (thaqāfa) — culture}

\label{lesson:AR-C267-thaqafa}



\begin{quote}
Before the new one: say the Arabic for a courtyard, then the Arabic for a tradition.
\end{quote}

\subsection*{You'll want to know: \ar{ثقافة}}
\textbf{\ar{ثقافة}} — \emph{thaqāfa} — \textquotedblleft{}culture\textquotedblright{}.

\begin{cousinweb}
\textbf{\ar{ثقافة}} (\emph{thaqāfa}) is built on the native Arabic root \emph{th-q-f}, \textquotedblleft{}to be skilful, to be sharp\textquotedblright{} (Lane). \emph{Thaqāfa} is skill, cultivation; the same root gives \emph{thiqāf}, a tool for straightening spears.
\end{cousinweb}

\begin{grammarlens}[title={What you've built}]
One more word for this chapter.
\end{grammarlens}

\subsection*{Your turn}
\begin{itemize}
\raggedright
  \item \emph{Say it:} \emph{thaqāfa}
  \item \emph{Say it:} \emph{thaqāfa}, once more
  \item \emph{Say it:} say \emph{thaqāfa}
  \item \emph{Recall:} say the Arabic for a courtyard, then the Arabic for a tradition, then say \emph{thaqāfa} again
\end{itemize}

\begin{tcolorbox}[breakable,colback=teal!4,colframe=teal!35!black,title={Before you move on}]
What does \ar{ثقافة} mean? (\textquotedblleft{}Culture\textquotedblright{}.) Say it once more.
\end{tcolorbox}
\section[\texorpdfstring{\ar{دين} (dīn)}{دين (dīn)}]{\ar{دين} (dīn) — religion}

\label{lesson:AR-C267-din}



\begin{quote}
Before the new one: say the Arabic for a tradition, then the Arabic for culture.
\end{quote}

\subsection*{You'll want to know: \ar{دين}}
\textbf{\ar{دين}} — \emph{dīn} — \textquotedblleft{}religion\textquotedblright{}.

\begin{cousinweb}
\textbf{\ar{دين}} (\emph{dīn}) is listed by Lane under \emph{d-y-n} with the senses recompense, obedience and religion. Some trace the sense \textquotedblleft{}religion\textquotedblright{} to Middle Persian \emph{dēn}, \textquotedblleft{}religion\textquotedblright{}; the question is not settled.
\end{cousinweb}

\begin{grammarlens}[title={What you've built}]
One more word for this chapter.
\end{grammarlens}

\subsection*{Your turn}
\begin{itemize}
\raggedright
  \item \emph{Say it:} \emph{dīn}
  \item \emph{Say it:} \emph{dīn}, once more
  \item \emph{Say it:} say \emph{dīn}
  \item \emph{Recall:} say the Arabic for a tradition, then the Arabic for culture, then say \emph{dīn} again
\end{itemize}

\begin{tcolorbox}[breakable,colback=teal!4,colframe=teal!35!black,title={Before you move on}]
What does \ar{دين} mean? (\textquotedblleft{}Religion\textquotedblright{}.) Say it once more.
\end{tcolorbox}
\section[\texorpdfstring{\ar{صلاة} (ṣalāh)}{صلاة (ṣalāh)}]{\ar{صلاة} (ṣalāh) — prayer}

\label{lesson:AR-C267-salah}



\begin{quote}
Before the new one: say the Arabic for culture, then the Arabic for religion.
\end{quote}

\subsection*{You'll want to know: \ar{صلاة}}
\textbf{\ar{صلاة}} — \emph{ṣalāh} — \textquotedblleft{}prayer\textquotedblright{}.

\begin{cousinweb}
\textbf{\ar{صلاة}} (\emph{ṣalāh}) is listed by Lane under \emph{ṣ-l-w}. It is usually traced to Aramaic \emph{ṣlōthā}, \textquotedblleft{}prayer\textquotedblright{}; Lane's sources say the related \emph{ṣalawāt}, \textquotedblleft{}synagogues\textquotedblright{}, is from Hebrew.
\end{cousinweb}

\begin{grammarlens}[title={What you've built}]
One more word for this chapter.
\end{grammarlens}

\subsection*{Your turn}
\begin{itemize}
\raggedright
  \item \emph{Say it:} \emph{ṣalāh}
  \item \emph{Say it:} \emph{ṣalāh}, once more
  \item \emph{Say it:} say \emph{ṣalāh}
  \item \emph{Recall:} say the Arabic for culture, then the Arabic for religion, then say \emph{ṣalāh} again
\end{itemize}

\begin{tcolorbox}[breakable,colback=teal!4,colframe=teal!35!black,title={Before you move on}]
What does \ar{صلاة} mean? (\textquotedblleft{}Prayer\textquotedblright{}.) Say it once more.
\end{tcolorbox}
\section[\texorpdfstring{\ar{علاقة} (ʿalāqa)}{علاقة (ʿalāqa)}]{\ar{علاقة} (ʿalāqa) — a relationship}

\label{lesson:AR-C267-alaqa}



\begin{quote}
Before the new one: say the Arabic for religion, then the Arabic for prayer.
\end{quote}

\subsection*{You'll want to know: \ar{علاقة}}
\textbf{\ar{علاقة}} — \emph{ʿalāqa} — \textquotedblleft{}a relationship\textquotedblright{}.

\begin{cousinweb}
\textbf{\ar{علاقة}} (\emph{ʿalāqa}) is built on the native Arabic root \emph{ʿ-l-q}, \textquotedblleft{}to cling, to hang\textquotedblright{} (Lane). \emph{ʿAlāqa} is a tie, a connection.
\end{cousinweb}

\begin{grammarlens}[title={What you've built}]
One more word for this chapter.
\end{grammarlens}

\subsection*{Your turn}
\begin{itemize}
\raggedright
  \item \emph{Say it:} \emph{ʿalāqa}
  \item \emph{Say it:} \emph{ʿalāqa}, once more
  \item \emph{Say it:} say \emph{ʿalāqa}
  \item \emph{Recall:} say the Arabic for religion, then the Arabic for prayer, then say \emph{ʿalāqa} again
\end{itemize}

\begin{tcolorbox}[breakable,colback=teal!4,colframe=teal!35!black,title={Before you move on}]
What does \ar{علاقة} mean? (\textquotedblleft{}A relationship\textquotedblright{}.) Say it once more.
\end{tcolorbox}
